%% Question bank chapter file - generated by bank_tools.py build (do not edit while the bank is still being built:
%% it is regenerated from the per-test data in Question Bank\data). Every question is one bankq block:
%%   \begin{bankq}[id=..., source={...}, paper=.., qno=.., type=..., subtopic={...}, answer={...}, answer_src=key|solved,
%%                 working={...}, difficulty=1-10, diff_parts={K S I C F}, diff_note={...},
%%                 group=<passage id or empty>, also_in={duplicate copies in other tests}]
%% difficulty = min(10, K+S+I+C+F-1): K concepts 1-3, S steps 1-4, I insight 0-2, C computation 0-2, F format 0-1
%%   ...question LaTeX, options, figures (\bfig / \bfigi -> img/)...
%%   \end{bankq}
%% Passages shared by several questions are bankpassage blocks; their questions carry group=<passage id>.
\documentclass[10pt]{article}
\usepackage[a4paper,margin=1.8cm]{geometry}
\usepackage[T1]{fontenc}
\usepackage{amsmath,amssymb,mathtools}
\usepackage{graphicx}
\usepackage[export]{adjustbox}
\usepackage{xcolor,needspace,keyval}
\usepackage{enumitem}
\usepackage{booktabs,array,tabularx}
\usepackage[hidelinks]{hyperref}
\setlength{\parindent}{0pt}
\setlength{\parskip}{3pt}
\makeatletter
\define@key{bq}{id}{\def\bq@id{#1}}\define@key{bq}{source}{\def\bq@src{#1}}\define@key{bq}{paper}{\def\bq@paper{#1}}
\define@key{bq}{qno}{\def\bq@qno{#1}}\define@key{bq}{type}{\def\bq@type{#1}}\define@key{bq}{subtopic}{\def\bq@sub{#1}}
\define@key{bq}{answer}{\def\bq@ans{#1}}\define@key{bq}{answer_src}{\def\bq@asrc{#1}}\define@key{bq}{working}{\def\bq@work{#1}}
\define@key{bq}{group}{\def\bq@group{#1}}\define@key{bq}{also_in}{\def\bq@also{#1}}
\define@key{bq}{difficulty}{\def\bq@diff{#1}}\define@key{bq}{diff_parts}{\def\bq@dparts{#1}}\define@key{bq}{diff_note}{\def\bq@dnote{#1}}
\newenvironment{bankq}[1][]{\def\bq@work{}\def\bq@also{}\def\bq@group{}\def\bq@diff{?}\def\bq@dparts{}\def\bq@dnote{}\setkeys{bq}{#1}%
  \par\needspace{4\baselineskip}\medskip\noindent{\small\textbf{\bq@id}\hfill\color{gray}[\bq@src, Paper \bq@paper, Q\bq@qno]
  \bq@type\ $\cdot$ \bq@sub\ $\cdot$ \textbf{D\bq@diff}}\par\nopagebreak\relax}
  {\par{\small\color{teal}Answer: \textbf{\bq@ans}\ifx\bq@work\@empty\else\ (solved: \bq@work)\fi
  \ifx\bq@also\@empty\else\par\color{gray}Also in: \bq@also\fi}\par}
\define@key{bp}{id}{\def\bp@id{#1}}\define@key{bp}{source}{\def\bp@src{#1}}\define@key{bp}{questions}{\def\bp@qs{#1}}
\newenvironment{bankpassage}[1][]{\setkeys{bp}{#1}\par\needspace{4\baselineskip}\bigskip\noindent
  {\small\textbf{Passage \bp@id}\hfill\color{gray}[\bp@src; questions \bp@qs]}\par\nopagebreak\relax\itshape}{\par}
\makeatother
\newcommand{\bfig}[2]{\par\smallskip{\centering\includegraphics[width=#1pt,max width=0.62\linewidth,max height=0.32\textheight]{img/#2}\par}\smallskip}
\newcommand{\bfigi}[2]{\includegraphics[width=#1pt,max width=\linewidth,max height=0.32\textheight,valign=c]{img/#2}}
\newcommand{\bmissing}[1]{\textit{[figure missing in the source conversion: \detokenize{#1}]}}
\newenvironment{qoptions}{\begin{itemize}[leftmargin=1.8em,itemsep=1pt,topsep=2pt,label={}]}{\end{itemize}}
\newcommand{\opt}[2]{\item[(#1)] #2}
\providecommand{\fwstart}{}\providecommand{\fwend}{}\providecommand{\qfullsplit}{}
\newcolumntype{L}[1]{>{\raggedright\arraybackslash}p{#1}}
\newcolumntype{C}[1]{>{\centering\arraybackslash}p{#1}}
\newcolumntype{R}[1]{>{\raggedleft\arraybackslash}p{#1}}
\providecommand{\sectionhdr}[1]{}\providecommand{\subsectionhdr}[1]{}


\begin{document}
{\Large\bfseries Solutions --- Question Bank}\par\smallskip
{\small Chemistry chapter C11. Questions from all converted test papers, grouped by source series. Answers: printed key where it matches the paper, otherwise solved (answer\_src=solved, with working). Exact repeats across papers appear once, with also\_in listing the other copies.}\par\medskip\hrule

\section*{Score-I papers}
\begin{bankq}[id=SI12-P2-CHEM-S3-Q09, source={Score-I - Test 12 (09 Mar 2022)}, paper=2, qno=9, type=matching, subtopic={Raoult's law}, answer={P$\to$2, Q$\to$1, R$\to$3}, answer_src=key, difficulty=2, diff_parts={K1 S1 I0 C0 F1}, diff_note={Recall azeotropes: hexane-heptane ideal, acetone-chloroform maximum, ethanol-water minimum}]
\begin{tabular}[t]{|p{0.65cm}|p{2.4cm}|p{0.65cm}|p{2.8cm}|}
\hline
 & \centering List-I & & \centering List-II\tabularnewline
\hline
(P) & n-hexane + n-heptane & (1) & Maximum boiling azeotrope\\
(Q) & Acetone + Chloroform & (2) & Can be separated by distillation\\
(R) & Ethanol + water & (3) & Minimum boiling azeotrope\\
\hline
\end{tabular}
\end{bankq}

\section*{TOAS papers}
\begin{bankq}[id=TOAS03-P1-CHEM-S1-Q03, source={TOAS - Test 03 (27 Feb 2022)}, paper=1, qno=3, type=single, subtopic={Colligative properties}, answer={A}, answer_src=key, difficulty=2, diff_parts={K1 S2 I0 C0 F0}, diff_note={van't Hoff factors 5, 3, 2 order the freezing point depressions}]
Assuming 100\% ionization, the increasing order of the freezing point of the solution will be :-
\begin{qoptions}
\opt{A}{$0.10\ \mathrm{mol\,kg^{-1}\ Ba_3(PO_4)_2}$ $<$ $0.10\ \mathrm{mol\,kg^{-1}\ Na_2SO_4}$ $<$ $0.10\ \mathrm{mol\,kg^{-1}\ KCl}$}
\opt{B}{$0.10\ \mathrm{mol\,kg^{-1}\ KCl}$ $<$ $0.10\ \mathrm{mol\,kg^{-1}\ Na_2SO_4}$ $<$ $0.10\ \mathrm{mol\,kg^{-1}\ Ba_3(PO_4)_2}$}
\opt{C}{$0.10\ \mathrm{mol\,kg^{-1}\ Na_2SO_4}$ $<$ $0.10\ \mathrm{mol\,kg^{-1}\ Ba_3(PO_4)_2}$ $<$ $0.10\ \mathrm{mol\,kg^{-1}\ KCl}$}
\opt{D}{$0.10\ \mathrm{mol\,kg^{-1}\ KCl}$ $<$ $0.10\ \mathrm{mol\,kg^{-1}\ Ba_3(PO_4)_2}$ $<$ $0.10\ \mathrm{mol\,kg^{-1}\ Na_2SO_4}$}
\end{qoptions}
\end{bankq}
\begin{bankq}[id=TOAS03-P1-CHEM-S2-Q09, source={TOAS - Test 03 (27 Feb 2022)}, paper=1, qno=9, type=multi, subtopic={Colligative properties}, answer={B,D}, answer_src=key, difficulty=4, diff_parts={K1 S3 I0 C0 F1}, diff_note={Compare osmotic pressures 0.2 vs 0.15 M and the direction of water flow}]
Study the following figure, and select the correct statements (assuming complete dissociation of the electrolyte).
\bfig{90.8}{TOAS03_p1_chem_s2_q09_1.png}
\begin{qoptions}
\opt{A}{There will be no movement of any solution across the membrane.}
\opt{B}{Water from $\mathrm{BaCl_2}$ will flow towards the NaCl solution.}
\opt{C}{Water from NaCl will flow towards the $\mathrm{BaCl_2}$.}
\opt{D}{The osmotic pressure of 0.1 M NaCl is higher than the osmotic pressure of 0.05 M $\mathrm{BaCl_2}$,}
\end{qoptions}
\end{bankq}
\begin{bankq}[id=TOAS03-P2-CHEM-S1-Q04, source={TOAS - Test 03 (27 Feb 2022)}, paper=2, qno=4, type=multi, subtopic={Raoult's law}, answer={A,B,C}, answer_src=key, difficulty=4, diff_parts={K1 S2 I1 C0 F1}, diff_note={P=P_E+(P_M-P_E)x gives P_E=135, P_M=254 and the equimolar value}]
At $40^\circ$C, vapour pressure in Torr of methanol and ethanol solution is $P=119x+135$ where $x$ is the mole fraction of methanol. Hence
\begin{qoptions}
\opt{A}{vapour pressure of pure methanol is 254 Torr}
\opt{B}{vapour pressure of pure ethanol is 135 Torr}
\opt{C}{vapour pressure of equimolar mixture of each is 194.5 Torr}
\opt{D}{vapour pressure of pure methanol is 119 Torr}
\end{qoptions}
\end{bankq}
\begin{bankq}[id=TOAS03-P2-CHEM-S2-Q03, source={TOAS - Test 03 (27 Feb 2022)}, paper=2, qno=3, type=numerical, subtopic={Colligative properties}, answer={374.15}, answer_src=key, difficulty=5, diff_parts={K2 S3 I0 C0 F1}, diff_note={Boiling-point elevation with NaCl (i=2) and sugar in 260 g water}]
The normal boiling point of pure water is 373.15 K. Calculate the boiling point (in K) of an aqueous solution containing 11.7 gm NaCl and 34.2 gm sugar ($\mathrm{C_{12}H_{22}O_{11}}$) in 260 gm water. For water, $K_b=0.52$ K--Kg/mol. [Atomic mass of Na = 23, Cl = 35.5]
\end{bankq}
\begin{bankq}[id=TOAS03-P2-CHEM-S2-Q04, source={TOAS - Test 03 (27 Feb 2022)}, paper=2, qno=4, type=numerical, subtopic={Raoult's law}, answer={2.45}, answer_src=solved, working={$x_{CO_2}=P/K_H=3.2\times10^5/1.6\times10^8=2\times10^{-3}$; 500 g water = 27.8 mol; $n_{CO_2}=0.0557$ mol = 2.45 g}, difficulty=5, diff_parts={K1 S3 I0 C1 F1}, diff_note={Henry's law mole fraction 2e-3 in 500 g water gives 2.45 g CO2}]
Henry's law constant for $\mathrm{CO_2}$ in water is $1.6\times10^8$ Pa at 298K. The mass of $\mathrm{CO_2}$ (in gm) is 500g of soda water when packed under 3.2 bar pressure at 298K.
\end{bankq}
\begin{bankq}[id=TOAS04-P2-CHEM-S1-Q04, source={TOAS - Test 04 (15 Mar 2022)}, paper=2, qno=4, type=multi, subtopic={Van't Hoff factor}, answer={A,D}, answer_src=key, difficulty=6, diff_parts={K2 S3 I1 C0 F1}, diff_note={Observed molar mass 122 from elevation in boiling point, hence no association in acetone}]
1.22 gm $\mathrm{C_6H_5COOH}$ is added into 100 gm of acetone, elevation in B.P. is found to be $0.17\,^\circ$C. If $K_b$ (acetone) = 1.7 K - Kg mol$^{-1}$ then select the incorrect options
\begin{qoptions}
\opt{A}{Observed molar mass of benzoic acid is 61 gm/mole.}
\opt{B}{Observed molar mass of benzoic acid is 122 gm/mole}
\opt{C}{Benzoic acid neither dissociate nor associate in acetone}
\opt{D}{Benzoic acid exist as dimer in acetone due to intermolecular H-bonding}
\end{qoptions}
\end{bankq}
\begin{bankpassage}[id=TOAS04-P2-CHEM-S2-G7-8, source={TOAS - Test 04 (15 Mar 2022)}, questions={7-8}]
2 moles of a volatile liquid (A) and 3 moles of another volatile liquid (B) are present in a solution at $25^\circ$C temperature in a closed container. If vapour pressure of pure A and pure B at $25^\circ$C are $P^0_A=300$ torr, $P^0_B=200$ torr respectively (above solution follows Raoult's law of partial vapour pressure)
\end{bankpassage}
\begin{bankq}[id=TOAS04-P2-CHEM-S2-Q07, source={TOAS - Test 04 (15 Mar 2022)}, paper=2, qno=7, type=single, subtopic={Raoult's law}, answer={A}, answer_src=key, difficulty=5, diff_parts={K2 S3 I1 C0 F0}, diff_note={Vapour composition from partial pressures, then pressure above the condensed distillate}, group=TOAS04-P2-CHEM-S2-G7-8]
Calculate vapour pressure of distillate of the solution at $25^\circ$C temperature.
\begin{qoptions}
\opt{A}{250 torr}
\opt{B}{240 torr}
\opt{C}{280 torr}
\opt{D}{270 torr}
\end{qoptions}
\end{bankq}
\begin{bankq}[id=TOAS04-P2-CHEM-S2-Q08, source={TOAS - Test 04 (15 Mar 2022)}, paper=2, qno=8, type=single, subtopic={Raoult's law}, answer={A}, answer_src=key, difficulty=3, diff_parts={K1 S3 I0 C0 F0}, diff_note={No azeotrope for this system: Raoult's law ideal solution, mole fraction in first vapour 0.5}, group=TOAS04-P2-CHEM-S2-G7-8]
Which of the following is incorrect ?
\begin{qoptions}
\opt{A}{Above solution form an Azeotropic mixture}
\opt{B}{Above solution can be separated by fractional distillation}
\opt{C}{Mole fraction of A in the 1st drop of vapour at $25^\circ$C is 0.5}
\opt{D}{Vapour pressure of solution is 240 torr at $25^\circ$C}
\end{qoptions}
\end{bankq}
\begin{bankq}[id=TOAS05-P1-CHEM-S3-Q04, source={TOAS - Test 05 (03 Apr 2022)}, paper=1, qno=4, type=numerical, subtopic={Colligative properties}, answer={9.60}, answer_src=key, difficulty=5, diff_parts={K1 S3 I1 C0 F1}, diff_note={Remaining 250 mL is 0.4 M; pressure = C R T = 9.6 atm}]
What should be the minimum pressure (in atm) applied on 0.1 M glucose solution, filled in semi-permiable membrane, so as to filter away 750 mL water, leaving 250 mL of glucose solution at 300 K.\\\relax
[Given : $R=0.08$ atm L$\mathrm{K^{-1}}$ mol$^{-1}$]
\end{bankq}
\begin{bankq}[id=TOAS05-P2-CHEM-S1-Q09, source={TOAS - Test 05 (03 Apr 2022)}, paper=2, qno=9, type=single, subtopic={Raoult's law}, answer={A}, answer_src=key, difficulty=2, diff_parts={K1 S2 I0 C0 F0}, diff_note={Vapour composition y_A = x_A P_A^0/(x_A P_A^0+x_B P_B^0)=0.1}]
Two liquids $A$ and $B$ have, $P_A^0:P_B^0=1:3$ at a certain temperature. Assume $A$ and $B$ form an ideal solution, and the ratio of moles of $A$ to $B$ in the liquid phase are $1:3$, then mole fraction of $A$ in vapour phase in equilibrium with the solution will be -
\begin{qoptions}
\opt{A}{0.1}
\opt{B}{0.2}
\opt{C}{0.5}
\opt{D}{1.0}
\end{qoptions}
\end{bankq}
\begin{bankq}[id=TOAS05-P2-CHEM-S2-Q05, source={TOAS - Test 05 (03 Apr 2022)}, paper=2, qno=5, type=numerical, subtopic={Colligative properties}, answer={384}, answer_src=key, difficulty=7, diff_parts={K2 S3 I1 C1 F1}, diff_note={Freezing at -4 C gives molality 2.15; remaining water, ice by difference}]
585.5 g of 0.5 molal aqueous solution of sucrose can be cooled up to $-4^\circ$C before freezing. The amount of ice (in g) that will be separated out is \underline{\hspace{1.2cm}}.\\
\relax[Given : $K_f(\mathrm{H_2O})=1.86$ K kg mol$^{-1}$]\\
\relax[Give your answer as nearest integer].
\end{bankq}
\begin{bankq}[id=TOAS06-P1-CHEM-S3-Q02, source={TOAS - Test 06 (17 Apr 2022)}, paper=1, qno=2, type=numerical, subtopic={Raoult's law}, answer={5.00}, answer_src=key, difficulty=5, diff_parts={K1 S3 I0 C1 F1}, diff_note={Henry's law: mass dissolved proportional to pressure and amount of water}]
At 300K, 40 ml of $\mathrm{O_3}$(g) dissolves in 100 g of water at 1.0 atm. What mass of ozone (in gram) dissolved in 1600 g of water at a pressure of 4.0 atm at 300 K ?\\
\relax[Report your answer nearest integer value.]
\end{bankq}
\begin{bankq}[id=TOAS07-P1-CHEM-S1-Q03, source={TOAS - Test 07 (01 May 2022)}, paper=1, qno=3, type=single, subtopic={Raoult's law}, answer={B}, answer_src=key, difficulty=6, diff_parts={K2 S3 I1 C1 F0}, diff_note={Two vapour-pressure equations in PA and PB with changed composition}]
An ideal binary solution contains 4 mol of A and 6 mol of B. Its vapour pressure is found to be 460 mmHg. To this solution, when 10 mol of B is added, its vapour pressure becomes 480 mmHg. The vapour pressures of pure A and pure B, respectively, are
\begin{qoptions}
\opt{A}{300 mmHg, 600 mmHg}
\opt{B}{400 mmHg and 500 mmHg}
\opt{C}{300 mmHg, 500 mmHg}
\opt{D}{400 mmHg and 600 mmHg}
\end{qoptions}
\end{bankq}
\begin{bankpassage}[id=TOAS07-P2-CHEM-S2-G9-10, source={TOAS - Test 07 (01 May 2022)}, questions={9-10}]
Two liquids A and B mix to form an ideal binary liquid solution. On adding solid non-volatile solute C to the solution, A starts polymerising into a hexamer soluble solid `$\mathrm{A_6}$' (non-electrolyte) following zero-order kinetics with a rate constant of reaction $=\left(\dfrac{10^{-1}}{6}\ \text{moles/min.}\right)$ If initially 10 moles of A and 20 moles of B were taken to form liquid solution and 5 moles of C were added, then answer the following questions based on above information and data given below :\\
\relax[Data : vapour pressure of pure A $=100$ torr, vapour pressure of pure B $=90$ torr].
\end{bankpassage}
\begin{bankq}[id=TOAS07-P2-CHEM-S2-Q09, source={TOAS - Test 07 (01 May 2022)}, paper=2, qno=9, type=single, subtopic={Raoult's law}, answer={A}, answer_src=key, difficulty=6, diff_parts={K2 S3 I1 C1 F0}, diff_note={A hexamerises (zero order): moles of A after 1 h, vapour pressure}, group=TOAS07-P2-CHEM-S2-G9-10]
The vapour pressure of the solution 1 hour after solid C is added will be given by :
\begin{qoptions}
\opt{A}{$\dfrac{220}{3}$ torr}
\opt{B}{98 torr}
\opt{C}{$\dfrac{2220}{29}$ torr}
\opt{D}{$\dfrac{2220}{24}$ torr}
\end{qoptions}
\end{bankq}
\begin{bankq}[id=TOAS07-P2-CHEM-S2-Q10, source={TOAS - Test 07 (01 May 2022)}, paper=2, qno=10, type=single, subtopic={Raoult's law}, answer={D}, answer_src=key, difficulty=6, diff_parts={K2 S3 I1 C1 F0}, diff_note={Time dependence of vapour pressure as A polymerises to completion}, group=TOAS07-P2-CHEM-S2-G9-10]
Which of the following statement in incorrect regarding the above process ?
\begin{qoptions}
\opt{A}{The vapour pressure will keep on decreasing initially on adding solute C.}
\opt{B}{The vapour pressure will become constant 100 min after addition of C.}
\opt{C}{Two hours after addition of C, the vapour pressure will become 67.5 torr.}
\opt{D}{Four hours after addition of C, the vapour pressure will become 73.75 torr.}
\end{qoptions}
\end{bankq}
\begin{bankq}[id=TOAS08-P1-CHEM-S1-Q04, source={TOAS - Test 08 (15 May 2022)}, paper=1, qno=4, type=multi, subtopic={Raoult's law}, answer={A,B}, answer_src=key, difficulty=7, diff_parts={K2 S4 I1 C0 F1}, diff_note={Equal moles of benzene and toluene: vapour pressure relations and the P-x graph}]
1 mole each of benzene and toluene are mixed (ideal solution) at a given temperature. A total of 1 mole has vaporised, at given pressure. Then, which of the following must be true for this equilibrium mixture? [Given : $P^0_{\text{Benzene}}=x$; $P^0_{\text{Toluene}}=y$]. Assume $X$, $U$, $V''$, $U'$, $Z$, $V'$, $Z'$ are various points on the curve.
\begin{qoptions}
\opt{A}{Total vapour pressure $<\left[\dfrac{x+y}{2}\right]$}
\opt{B}{Total vapour pressure $=\sqrt{xy}$}
\opt{C}{In the graph,\newline \bfigi{117.8}{TOAS08_p1_chem_s1_q04_1.png}\newline the vapour pressure of equilibrium mixture may correspond to point $V''$.}
\opt{D}{If in the graph given in option(C), external pressure is isothermally increased, then last trace of liquid will disappear at point $Z$.}
\end{qoptions}
\end{bankq}
\begin{bankq}[id=TOAS08-P1-CHEM-S3-Q05, source={TOAS - Test 08 (15 May 2022)}, paper=1, qno=5, type=numerical, subtopic={Colligative properties}, answer={3.00}, answer_src=key, difficulty=7, diff_parts={K2 S3 I1 C1 F1}, diff_note={Elevation in boiling point of a saturated solution gives the van't Hoff factor}]
The saturated solution of a certain salt show elevation in boiling point equal to $3.12\times10^{-3}$ K in aqueous solution. If the solubility of salt is $2\times10^{-3}$ mole/kg at boiling point, then how many moles of ions are given by 1 mole of solid salt on complete dissociation.\\
$K_{b(\mathrm{H_2O})}=0.52$ K kg/mole
\end{bankq}
\begin{bankq}[id=TOAS08-P2-CHEM-S1-Q03, source={TOAS - Test 08 (15 May 2022)}, paper=2, qno=3, type=multi, subtopic={Raoult's law}, answer={B,D}, answer_src=key, difficulty=4, diff_parts={K1 S3 I0 C0 F1}, diff_note={Negative deviation pairs: hydrogen bonding and dipole attractions}]
Identify which of the following solution(s) will show negative deviation from Raoult's law -
\begin{qoptions}
\opt{A}{Carbon-disulphide $+$ acetic acid}
\opt{B}{Water $+$ Formic Acid}
\opt{C}{Propan-2-ol $+$ 2-methyl propanol}
\opt{D}{Chloroform $+$ Acetone}
\end{qoptions}
\end{bankq}
\begin{bankq}[id=TOAS08-P2-CHEM-S1-Q04, source={TOAS - Test 08 (15 May 2022)}, paper=2, qno=4, type=multi, subtopic={Raoult's law}, answer={A,B,C}, answer_src=key, difficulty=7, diff_parts={K2 S4 I1 C0 F1}, diff_note={Benzene-toluene solution: vapour pressure and vapour composition}]
Pure benzene has a vapour pressure 900 torr and pure toluene has a vapour pressure 360 torr. If a solution is prepared by mixing 39gm benzene and 46 gm toluene then
\begin{qoptions}
\opt{A}{Vapour pressure of its distillate will be be 745.714 torr}
\opt{B}{Vapour pressure of initial solution will be 630 torr}
\opt{C}{Mole \% of benzene in vapour phase of solution will be $\dfrac{500}{7}\%$}
\opt{D}{Mole \% of toluene in its distillate will be 50\%.}
\end{qoptions}
\end{bankq}
\begin{bankq}[id=TOAS09-P1-CHEM-S2-Q10, source={TOAS - Test 09 (30 May 2022)}, paper=1, qno=10, type=multi, subtopic={Raoult's law}, answer={A,B,C,D}, answer_src=key, difficulty=6, diff_parts={K2 S3 I1 C0 F1}, diff_note={Raoult's law: standard approach for this subtopic; 2 concept(s), 3 step(s), needs a non-obvious insight}]
Two volatile liquids A \& B are mixed to form an ideal solution. If $P_A^0=100$ mm of Hg \& $P_B^0=150$ mm of Hg then identify the correct option(s) if 10 moles of A are mixed with 40 moles of B such that moles vaporised are very small.
\begin{qoptions}
\opt{A}{Total vapour pressure will be 140 mm of Hg.}
\opt{B}{Mole fraction of A in vapour state will be less than 0.5.}
\opt{C}{Mole fraction of B in vapour state is equal to $\dfrac{6}{7}$}
\opt{D}{Mole \% of B in vapour state will be greater than 80\%.}
\end{qoptions}
\end{bankq}
\begin{bankq}[id=TOAS09-P2-CHEM-S1-Q14, source={TOAS - Test 09 (30 May 2022)}, paper=2, qno=14, type=single, subtopic={Raoult's law}, answer={C}, answer_src=key, difficulty=4, diff_parts={K2 S2 I1 C0 F0}, diff_note={Raoult's law: standard approach for this subtopic; 2 concept(s), 2 step(s), needs a non-obvious insight}]
At 300 K, the vapour pressure of an ideal solution containing 3 mole of A and 2 mole of B is 600 torr. At the same temperature, if 1.5 mole of A \& 0.5 mole of C (non-volatile) are added to this solution the vapour pressure of solution increases by 30 torr. What is the value of $P_B^{\circ}$?
\begin{qoptions}
\opt{A}{940 torr}
\opt{B}{405 torr}
\opt{C}{90 torr}
\opt{D}{100 torr}
\end{qoptions}
\end{bankq}
\begin{bankq}[id=TOAS09-P2-CHEM-S1-Q15, source={TOAS - Test 09 (30 May 2022)}, paper=2, qno=15, type=single, subtopic={Van't Hoff factor}, answer={D}, answer_src=key, difficulty=4, diff_parts={K2 S2 I1 C0 F0}, diff_note={Van't Hoff factor: standard approach for this subtopic; 2 concept(s), 2 step(s), needs a non-obvious insight}]
What is the correct sequence of osmotic pressures of 0.01 M aq. solutions of :\par
(i) Al$_2$(SO$_4$)$_3$ ($\pi_1$)\par
(ii) Na$_3$PO$_4$ ($\pi_2$)\par
(iii) BaCl$_2$ ($\pi_3$)\par
(iv) Glucose ($\pi_4$)\par
[Assume 100\% ionisation of electrolytes]
\begin{qoptions}
\opt{A}{$\pi_4>\pi_2>\pi_3>\pi_1$}
\opt{B}{$\pi_3>\pi_4>\pi_2>\pi_1$}
\opt{C}{$\pi_3>\pi_4>\pi_1>\pi_2$}
\opt{D}{$\pi_1>\pi_2>\pi_3>\pi_4$}
\end{qoptions}
\end{bankq}
\begin{bankq}[id=TOAS09-P2-CHEM-S1-Q16, source={TOAS - Test 09 (30 May 2022)}, paper=2, qno=16, type=single, subtopic={Raoult's law}, answer={C}, answer_src=key, difficulty=4, diff_parts={K2 S2 I1 C0 F0}, diff_note={Raoult's law: standard approach for this subtopic; 2 concept(s), 2 step(s), needs a non-obvious insight}]
Find ratio of mole fractions of gases N$_2$ \& O$_2$ dissolved in water from air at 25$^{\circ}$ \& 1 atm. Assume air is 20\,\% by volume of O$_2$ \& 80\% by volume of N$_2$\par
\textbf{Given :} $K_H(\mathrm{N_2})=2\times10^4$ atm\par
\hspace*{3.3em}$K_H(\mathrm{O_2})=10^4$ atm
\begin{qoptions}
\opt{A}{8 : 1}
\opt{B}{1 : 8}
\opt{C}{2 : 1}
\opt{D}{1 : 1}
\end{qoptions}
\end{bankq}
\end{document}
